\documentclass[11pt,a4paper]{article}
\usepackage[margin=23mm]{geometry}
\usepackage{amsmath,amssymb,booktabs,array}
\input{glyphtounicode}
\pdfgentounicode=1
\usepackage[hidelinks]{hyperref}
\hypersetup{hypertexnames=false,pdftitle={From S13 to the current VSL model: a unified recap}}
\setlength{\parindent}{0pt}
\setlength{\parskip}{6pt}
\setlength{\emergencystretch}{2em}
\newcommand{\pos}[1]{[#1]_+}
\newcommand{\budget}{\mathcal B}
\newcommand{\loss}{p_{\mathrm{rec}}}
\newcommand{\stackbox}[1]{\fbox{\parbox{66mm}{\centering #1}}}
\renewcommand{\labelitemi}{\(\bullet\)}
\begin{document}
\begin{center}
{\LARGE From S13 to the current VSL model}\\[5pt]
{\large One traffic chain, with the changes made explicit}\\[5pt]
Feedback 8 recap --- 16 October 2026
\end{center}

\section*{1. The backbone that remains}
The task is still a permanent three-to-two lane drop with one finite traffic
cohort and one constant speed advice shared by all lanes. The research asks
whether information about unfinished mandatory merging helps predict a useful
speed decision. It does not assume that a full lane representation, vehicle
identity or discrete order is necessary.

S13's central idea was simple: slower travel gives traffic more time to
prepare for the lane drop; better preparation can reduce the work required
at the final merge. The total-time objective charges the slower travel as
well as subsequent service.

The current stack keeps the control question while changing the traffic
mechanism: progress and spare receiving govern early transfer; timed
arrivals feed a coupled choke; admitted E can restrict later service.
The two backbones can therefore be read side by side:
\begin{center}
\begin{minipage}[t]{0.48\textwidth}
\centering\small
\textbf{S13: preparation time and batch work}\\[7pt]
\stackbox{Initial ending / continuing mass\\Common speed advice}\\[3pt]
$\downarrow$\\[3pt]
\stackbox{Preparation time\\$\to$ conservative early transfer}\\[3pt]
$\downarrow$\\[3pt]
\stackbox{Remaining ending mass\\$\to$ fixed receiving work}\\[3pt]
$\downarrow$\\[3pt]
\stackbox{Proportional batch service\\$\to$ uniform departures}\\[3pt]
$\downarrow$\\[3pt]
\stackbox{Complete cohort time\\$\to$ analytical sensitivity}
\end{minipage}\hfill
\begin{minipage}[t]{0.48\textwidth}
\centering\small
\textbf{Current: traffic state and shared service}\\[7pt]
\stackbox{Initial footprint and lane masses\\Common speed advice}\\[3pt]
$\downarrow$\\[3pt]
\stackbox{Progress-based E $\to$ A merging\\Limited by A's spare receiving}\\[3pt]
$\downarrow$\\[3pt]
\stackbox{Fixed cutoff; finite terminal travel\\$\to$ staggered eligible arrivals}\\[3pt]
$\downarrow$\\[3pt]
\stackbox{\textbf{Coupled point merge}\\A/E queues, speed marks and work\\[2pt]
\begin{tabular}{@{}c@{\,}c@{\,}c@{}}
Service & $\xrightarrow{\text{E entries}}$ & recovery\\
Service & $\xleftarrow{\text{available budget}}$ & recovery
\end{tabular}}\\[3pt]
$\downarrow$\\[3pt]
\stackbox{Departures through full clearance\\$\to$ complete local time}
\end{minipage}
\end{center}
E is the ending lane; A is its adjacent continuing lane. In the current
choke, E entries update recovery, which changes the budget for later
service. The arrows inside that box form a causal feedback loop.

\textbf{Reading the status correctly.} This is a functionally complete
\emph{conditional mathematical model}. The selected formulas and their
forward arguments are assembled here; this recap changes none of them.
Physical parameter admission, the amended information/action and gradient
protocol, and the final SUMO correspondence remain unfinished. Step 3 stays
open at MP6-Q1; Step 4 is unopened. No implementation or experiment is reported.

\newpage
\section*{2. What changed: setting, motion and early merging}
The component numbers preserve the original nine-component audit. Several
equations belong to the same component; they are not additional project stages.

\begin{tabular}{>{\raggedright\arraybackslash}p{31mm}>{\raggedright\arraybackslash}p{119mm}}
\toprule
Original component & Old $\to$ current; the core question it now answers\\
\midrule
\textbf{1. Geometry, population and preparation} &
\textbf{Modified.} S13 grouped both continuing lanes and placed the cohort
at one preparation entrance. The current model distinguishes ending E,
adjacent A and far continuing B, with a uniform occupied footprint and an
initially empty terminal. \emph{Which lane receives the merge, and where is
the traffic when advice begins?} Initial A/B symmetry remains; their later
states need not be equal.\\[7pt]
\textbf{2. Advice and response} &
\textbf{Retained with an explicit boundary.} One bounded, constant common
advice remains. It sets the ideal free-motion speed in the admitted branch;
queued boundary marks can be lower. \emph{Through which physical state can
advice influence service?} The final work law has fixed coefficients and
no direct advice multiplier. Instantaneous native compliance is unverified.\\[7pt]
\textbf{3. Longitudinal timing} &
\textbf{Replaced.} One shared traversal time followed by batch service becomes
position-dependent free travel, a fixed early-merge cutoff, and finite
terminal travel/storage represented by boundary histories. \emph{When does
each part of the cohort become available at the choke, and can its queue
reach preparation?} No new spatial grid is introduced.\\[7pt]
\textbf{4. Mandatory preparation transfer} &
\textbf{Replaced.} A constant completion fraction per second becomes a
homogeneous offer limited by donor traffic and A's spare receiving.
\emph{How much early merging can actually succeed?} A fixed physical
opportunity length replaces an unqualified elapsed-time completion law.\\[7pt]
\textbf{5. Preparation receiving} &
\textbf{Made explicit.} S13 assumed it did not bind. A's own traffic now
claims receiving first; E uses the remaining work opportunity. Conservative
terminal storage bounds protect this homogeneous upstream calculation.
\emph{Can A accommodate the attempted transfers without an omitted
upstream queue?}\\
\bottomrule
\end{tabular}

Components 4 and 5 form one joined preparation component in the present
equations. They are listed separately to show which original assumptions
changed. Distinguishing A from B does not create new initial information
requirements under the retained initial symmetry.

\newpage
\section*{3. What changed: final merging, accounting and evidence}
\begin{tabular}{>{\raggedright\arraybackslash}p{31mm}>{\raggedright\arraybackslash}p{119mm}}
\toprule
Original component & Old $\to$ current; the core question it now answers\\
\midrule
\textbf{6. Choke service and departures} &
\textbf{Replaced and expanded.} Fixed receiving work and proportional
batch service become a causal A/E node. It uses eligible queues, free/queued
speed marks, current maneuver work, a priority/cooperation rule and recent
E recovery history. \emph{Who can cross now, and how does an admitted slow
merge affect later crossing?} A can now wait as well as E. B remains
independent in the admitted domain.\\[8pt]
\textbf{7. Complete-time objective} &
\textbf{Meaning retained; calculation replaced.} Full affected-cohort time
remains the objective, but the old batch formula is replaced by actual
departure or queue-area integration. \emph{Does every unit's travel and
waiting count exactly once?} Completion is at the choke; downstream
physical influence may remain after that unit's own clock stops.\\[8pt]
\textbf{8. Information and initialization} &
\textbf{Retained and rechecked for the forward law.} Aggregate history,
that history plus initial E mass, and equivalent explicit lane state remain
the comparison ladder. \emph{What extra initial measurement is sufficient?}
Known total mass plus initial E mass still reconstructs this member.
Future queues and recovery history are predicted internally.\\[8pt]
\textbf{9. Sensitivity and decision evidence} &
\textbf{Standard retained; result must be requalified.} The physical common
advice is still the differentiated variable. \emph{Does the complete
objective's response support a useful decision?} S13's analytical action
split remains valid for S13, but does not automatically transfer to these
changed laws. The current stack may have nonsmooth activation events.\\
\bottomrule
\end{tabular}

The genuinely new dynamic content is concentrated in component 6: A's
waiting balance, the speed/work coupling, and the retained recovery account.
Finite terminal histories also add timing and storage information to
component 3. None of these is a microscopic gap or vehicle-identity model.

SUMO is an external evidence layer, not a component needed to calculate
the primary forward trajectory. The agreed ladder is primary-model
verification/optimization, restricted SUMO S-P, then natural late-merging
SUMO S-L transfer. Its amended protocol remains to be completed.

\newpage
\section*{4. The new preparation law: opportunity plus spare receiving}
Use E for the ending lane, A for its adjacent continuing lane and B for
the far lane. Let $e(t),a(t),b_0$ be densities in the homogeneous upstream
body. Put $r=e_0+a_0$, with $a_0=b_0$, and let $h>0$ be a fixed physical
opportunity length. Define
\[
 K=wk_j,\qquad k_c(v)=\frac{K}{w+v},\qquad C(v)=v k_c(v).
\]
Here $w$ is a backward-wave speed and $k_j$ jam density. These define the
selected triangular traffic relation; they are not fitted in this recap.

\textbf{The core rule.} A uses its own receiving first. E's accepted
opportunity-equivalent rate $f$ is the smaller of its offer and the spare
receiving available after that claim:
\[
 S(v,a)=\min\{C(v),w(k_j-a)\},\qquad
 f=\min\left\{ve,\frac{S(v,a)-va}{\gamma}\right\},\qquad
 \dot e=-f/h,\quad\dot a=f/h.
\]
$\gamma\geq1$ is early-merge receiving work. Equal mass leaves E and enters
A. $f$ is not an extra longitudinal boundary flow: the actual lateral
source per unit road length is $f/h$.

In the admitted free upstream branch, $S=C$ and $e+a=r$. Expressed per
distance traveled, $z=vt$, the same rule is
\[
 \frac{de}{dz}=-\frac1h\min\left\{e,
          \frac{\Delta(v)+e}{\gamma}\right\},\qquad
 \Delta(v)=k_c(v)-r.
\]
This exposes an important change from S13. If donor availability alone
limits transfer, the residual is exponential in \emph{distance}, so slower
motion alone gives no extra preparation over the same distance. A speed
effect survives through the declared receiving headroom when that limit
matters. Longer residence by itself is no longer assumed to guarantee
better preparation.

\textbf{How this becomes arrivals.} The preparation ends at $L_p$, before
a terminal of length $b$, with $L=L_p+b$. Initial positions occupy
$[0,\ell]$. Over the handoff interval
$((L_p-\ell)/v,L_p/v)$,
\[
 \dot F_E=ve(t),\quad \dot F_A=v(r-e(t)),\quad\dot F_B=vb_0;
 \qquad P_j(t)=F_j(t-b/v).
\]
$F_j$ counts terminal entries and $P_j$ counts traffic that has had time
to reach the choke. Early merging stops at $L_p$; remaining E must merge
at the choke. This is the selected opportunity restriction, not a claim
that native drivers universally avoid late lateral opportunities.

\newpage
\section*{5. The new point service: traffic states determine current work}
The terminal is a finite physical approach, represented by cumulative
histories rather than interior cells. For each lane,
\[
 N_j=F_j-Y_j,\qquad q_j=P_j-Y_j,\qquad g_j=\dot Y_j.
\]
$N_j$ is traffic physically in the terminal. $q_j$ is traffic eligible
to cross but still unserved. They are different quantities. Queues can
have physical extent; a common service point does not make them vertical
stacks or imply a stationary queue tail.

\textbf{Speed marks.} Free traffic uses $u_j=v$. A queued approach uses
the equilibrium boundary mark
\[
 u_j=U(g_j),\qquad U(g)=\frac{wg}{K-g}.
\]
This is a representative crossing-speed assumption, not a tracked vehicle
speed or a whole-lane average. Its physical accuracy remains to be tested.

\textbf{Immediate work.} With the available shared account $\budget(t)$,
\[
 c_A=1,\qquad c_E=c_0+\chi g_A\pos{u_A-u_E},\qquad
 c_0=1+\delta\geq1,\qquad g_A+c_Eg_E\leq\budget.
\]
$c_E$ is work per E mass. It includes a baseline maneuver burden and extra
work when passing A is faster than E. The coefficient $\chi$ is fixed:
advice influences this expression through traffic state, not through a
moving normalizer. There is no mismatch burden without passing A traffic.

\textbf{Who gets the service?} If A is free and its offer fits, it receives
its offer first. E takes feasible residual work. If A queues, cooperation
gives each backlogged stream equal work; unused work goes to a stream that
can use it. Equal work need not mean equal vehicle flow. For queued E,
define the unique service root by
\[
 H(g;a,u)=g\{c_0+\chi a\pos{u-U(g)}\}=x.
\]
Free A uses $(a,u,x)=(\lambda_A,v,\budget-\lambda_A)$.
Both queued use $(a,u,x)=(\budget/2,U(\budget/2),\budget/2)$.
Here $\lambda_j=\dot P_j$.

For example, free fast A beside queued slow E produces a large mismatch
burden and scarce E service. If A also queues, both marks are lower and
cooperation becomes active. A retains an advantage from E's maneuver
cost. Cooperation remains applicable while A queues even if E clears
first; it is not a literal two-positive-queues toggle.

An empty approach stays free while its offer fits. A strict service
shortfall starts waiting; true clearance retests free admission. Equality
alone does not create a queue, and no numerical queue threshold is added.

\newpage
\section*{6. The new recovery account: one reading, sufficient history}
Current entry and its later influence are two parts of the selected point
model. After a slow E entry, following traffic can remain restricted even
though that E traffic has crossed the completion line. The model retains
this influence in one shared account:
\[
 \budget(t)=\pos{R-\loss(t)}.
\]
$R$ is undisturbed shared service. The current model derives a strictly
positive lower budget bound in its conditional domain; it does not insert
an artificial service floor.

\textbf{A single entry's remaining influence.} With physical acceleration
parameter $a_r>0$, an admitted speed mark $u$ generates
\[
 V(\theta;u)=\sqrt{(w+u)^2+2wa_r\theta}-w,\qquad
 D(\theta;u)=\pos{R-c(V(\theta;u))},\quad c(s)=\frac{Ks}{w+s}.
\]
The profile decreases as it ages and becomes zero after a finite time.
$V$ is a transformed boundary quantity used to calculate passing capacity;
it is \emph{not} the vehicle's speed trajectory. $u$, by contrast, is the
model's representative insertion-speed mark.

\textbf{How entries combine.} Only E admissions create new profiles.
Let $T(\xi)$ be the time E mass rank $\xi$ crossed. With fixed physical
normalization $m=1$ vehicle-equivalent,
\[
 \loss(t)=\frac1m\int_{[Y_E(t)-m]_+}^{Y_E(t)}
 D(t-T(\xi);u_E(T(\xi)))\,d\xi.
\]
New E mass adds its profile and retires equal oldest weight once the
recent-mass window fills. Retained profiles continue aging. With no new
E, recovery continues. There is no per-step reset, queue-clearance reset,
or deletion of traffic when an old profile retires.

This makes the service reading a singleton, but its future generally
requires recent age/mark information. Two equal readings may recover
differently. That memory is generated by the model and is not extra future
information supplied to the controller.

The approximation can retire an old strong restriction when newer fast E
passes. A does not independently deposit a second restriction merely
because it is slowed by the existing one. Both are selected simplifications
whose physical credibility remains open. They are not exact microscopic
leader or gap dynamics.

\newpage
\section*{7. What is counted, and what information initializes it?}
\textbf{The objective remains complete cohort time.} Every initial unit
is counted until it crosses the choke. With total mass $M$,
\[
 J(v)=\int_0^\infty\left(M-\sum_{j=E,A,B}Y_j(t)\right)dt
 =\frac{M(L-\ell/2)}v+\int_0^\infty(q_A+q_E)\,dt.
\]
The first term is travel from the uniform initial footprint. The second
counts all excess terminal waiting, including A's new delay. B stays
unqueued under its adequate independent service condition. Dividing by
$M$ gives mean time and preserves action rankings.

The old S13 formula used half the batch clearance duration because its
departures were uniform over that duration. The current model cannot
replace queue area by half of final clearance: arrivals, priority and
recovery now determine the departure curve.

\textbf{Completion and physical influence are different.} Once an E unit
crosses, its own local time is finished. Its retained recovery profile can
still hold back unfinished A/E traffic. The waiting of those followers is
counted. Post-line travel time of the completed unit is outside the chosen
local objective. Recovery memory is not traffic mass.

\textbf{One extra initial scalar is still sufficient.} Let $E_0$ be
initial ending mass. Known $M,\ell$ and initial A/B symmetry give
\[
 e_0=E_0/\ell,\qquad a_0=b_0=(M-E_0)/(2\ell).
\]
Known empty terminal and recovery histories complete initialization. The
model then generates its own preparation, arrivals, queues and recovery.
Aggregate-plus-$E_0$ and explicit lane state therefore make the same
forward prediction for this member. Full lane state is not established
as a necessary remedy; the current meter reading alone is not sufficient
evolving state.

\textbf{Why the speed tradeoff is still an open result.} Speed affects
free travel, preparation headroom, arrival times and potentially free-A
entry mismatch. Recovery couples later service to earlier E admissions.
These paths can oppose one another. A slower optimum or useful gradient
does not follow merely because all paths are present. The earlier S13
action-split certificate belongs to its original equations.

\newpage
\section*{8. What has been established, and at which level?}
\begin{tabular}{>{\raggedright\arraybackslash}p{43mm}>{\raggedright\arraybackslash}p{107mm}}
\toprule
Question & Current answer\\
\midrule
Is there a complete forward law? & Yes conditionally: inputs and predicted
history determine preparation, entry, service, queues, complete departures
and TST throughout the original advice interval under the retained bounds.\\[6pt]
Can merging delay A? & Yes in a constructed complete analytical family,
with positive recovery and mismatch, ordinary underload and active renewal.
This is not a calibrated parameter range or a preferred-action result.\\[6pt]
Can storage protect homogeneous preparation? & Yes under the sufficient
finite-storage and entry-capacity conditions. The model does not thereby
solve arbitrary upstream spillback.\\[6pt]
Is extra initial information compact? & Yes. Total mass plus initial E mass
and symmetry reconstruct this member. A richer representation has no
promised advantage at the same information.\\[6pt]
Is the joined physical regime admitted? & No. Crossing-speed interpretation,
preparation scales, current work and recovery need joint evidence with one
parameter set. A good TST fit alone cannot qualify their separation.\\[6pt]
Are useful actions and gradients proved for this stack? & No. Their amended
review and independent evaluation remain. Forward uniqueness does not imply
smooth sensitivity or an interior optimum.\\
\bottomrule
\end{tabular}

The forward proof exploits the monotone H1 arrival pulse and selected
recovery/priority rules. E can have at most one waiting episode in that
domain; repeated A events have a positive separation. A finite cohort
\emph{alone} does not imply that E cannot requeue. Bursty later E arrivals
can requeue in a wider setting or native SUMO, which must be interpreted
as possible model discrepancy rather than automatically rejected.

A useful negative control is already established: remove the baseline
extra maneuver work ($\delta=0$), and ordinary underload keeps the primary
cohort free. Mismatch and recovery never activate, so fastest advice is
common. Positive information or optimizer results are not required.

\newpage
\section*{9. Literature models that are serious rivals}
The review has identified the following comparison candidates. They are
established model families with relevant mechanisms, not yet a frozen set
of benchmarks for this exact observation, action and objective contract.

\textbf{Aggregate history and capacity drop ---}
\href{https://doi.org/10.1111/mice.13492}{Wei, Han and Wang (2025)}.
A first-order link model retains cumulative counts and entry-speed history,
with VSL effects and traffic-state-dependent discharge loss. This is a close
aggregate rival: useful memory and endogenous service need not require
lane decomposition. It is stronger than a density-only, fixed-capacity peer.

\textbf{Compact lane-drop capacity control ---}
\href{https://arxiv.org/abs/1310.2658}{Jin and Jin}.
An aggregate lane-drop model represents a discontinuous reduction in
discharge and studies VSL control. It offers a simpler explanation of a
speed benefit through avoiding capacity loss. Its idealized limits also
help identify conditions where VSL has no benefit. Its capacity-drop
branch needs its own treatment before gradient comparisons.

\textbf{Dynamic aggregate speed and predictive control ---}
\href{https://www.dcsc.tudelft.nl/~bdeschutter/pub/rep/03_016.pdf}{Hegyi,
De Schutter and Hellendoorn (2005)}.
METANET predicts segment density and mean-speed dynamics under speed limits;
a predictive controller uses these to manage congestion. It provides an
aggregate control rival with evolving speed rather than an instantaneous
free-motion response. Its repeated decisions differ from our constant advice.

\textbf{Lane-aware receiving and priority ---}
\href{https://doi.org/10.1007/s11067-013-9193-7}{Carey, Balijepalli and Watling}.
A multi-lane CTM couples longitudinal and lateral transfers through finite
target receiving and extra lane-changing headway, with alternative priority
rules. It is both a source lineage and a fuller continuous lane-aware rival.
The published model does not itself supply our common-VSL comparison.

\textbf{Density-driven mandatory merging ---}
\href{https://doi.org/10.1080/21680566.2019.1700846}{Nagalur Subraveti,
Knoop and van Arem}.
Lane densities, lane-end incentives and receiving constraints determine
realized transfers. This is a spatial continuous alternative to homogeneous
preparation. Its inspected validation concerns lane-density prediction;
common-advice response and gradient usefulness require separate qualification.

\textbf{Analytical speed-control gradients ---}
\href{https://arxiv.org/abs/1603.04785}{Delle Monache, Piccoli and Rossi}.
An LWR speed-control formulation derives cost variations for outflow tracking.
It is a gradient-method reference, rather than a matched lane-drop task:
its control and objective differ. VSL gradients themselves are established.

\textbf{Comparison principle.} Published models test alternative traffic
explanations. Aggregate-plus-initial-E tests information sufficiency, while
bounded scalar search and finite differences test optimization methods.
Keep these roles distinct and match the admitted inputs, action and full
time account before attributing a performance difference to information.

\newpage
\section*{10. How to use the two reports}
This recap answers the framework and old-to-new questions. The companion
report supplies the equations, units, admissible domain and derivations
in one notation, including conservation, root uniqueness, recovery,
event continuation, full clearance, the complete analytical A-delay
construction and the physical-identification limits.

\begin{center}
\href{vsl_feedback8_full_mathematics.pdf}{\textbf{Companion: consolidated mathematical report}}
\end{center}

The mathematical report distinguishes the historical S13 certificate from
the selected current law. Historical rejected alternatives are not silently
combined into a primary model. It records $\budget$ for the shared budget
and $\loss$ for its restriction; older reports often write $B$ and $p$.
This is a notation clarification only, avoiding confusion with far lane B.

\textbf{Source basis and adaptation.} The component chain draws on the
already reviewed lane-conservation/mandatory-transfer literature, receiving-
work and allocation literature, boundary-count traffic propagation and
point-insertion recovery. H1's homogeneous law, the selected speed/work
closure, priority guards and recent-mass renewal are declared adaptations.
No cited paper supplies this complete joined stack. The detailed source
and proof map is retained separately:

\begin{center}
\href{../../extraction/research_readiness/vsl_feedback8_model_stack_source_map.md}{Archived source, equation and decision map}
\end{center}

\textbf{For an eventual control study.} Compare full aggregate history,
that history plus initial ending mass, and explicit lane state at the same
information. Keep capable aggregate, history and capacity-drop models,
analytical and finite-difference gradients, and bounded scalar search as
serious rivals. Correct differentiation, improved internal cost and external
policy usefulness remain different claims.

\textbf{Immediate disposition.} Feedback 8 is a consolidation request, not
authorization to resume the unfinished protocol reviews. This delivery
finishes the recap and full mathematical reporting. The active scientific
checkpoint remains MP6-Q1; applied MP6 and Step 3 stay open. After review,
the next allowed review work is the information/action/full-objective
gradient protocol and then the updated model $\to$ SUMO S-P $\to$ SUMO S-L
contract. Step 4, simulator development, calibration and traffic experiments
remain unopened. S13, the urban task and the ramp line retain their closure.
\end{document}
